\section{The killed functional}\label{sec:dv}

Every rate in this paper is the restriction of one functional, the Donsker and
Varadhan functional of the chain.

\subsection{The restriction identity}\label{sec:restriction}

For a probability vector $\alpha$ on $[d]$ let
\begin{equation}\label{eq:JV}
 J_V(\alpha)=\sup_{u\in(0,\infty)^d}\sum_{i\in[d]}\alpha_i\Bigl(-\frac{(Qu)_i}{u_i}\Bigr)
 =\sup_{u\in(0,\infty)^d}\sum_{i\in[d]}\alpha_i\Bigl(q_i-\sum_{j\ne i}q_{ij}\frac{u_j}{u_i}\Bigr),
\end{equation}
the Donsker and Varadhan rate of the empirical measure of the continuous-time
chain with generator $Q$~\cite{dv1975}. For a face $F$ and $\alpha\in\Delta_F$ let
\[
 J_F(\alpha)=\sup_{u\in(0,\infty)^F}\sum_{i\in F}\alpha_i\Bigl(-\frac{(Q_Fu)_i}{u_i}\Bigr)
 =\sup_{u\in(0,\infty)^F}\sum_{i\in F}\alpha_i\Bigl(q_i-\sum_{j\in F\setminus\{i\}}q_{ij}\frac{u_j}{u_i}\Bigr),
\]
the same functional for the chain killed on leaving $F$. Let $q_i^{\rm out}(F)=\sum_{j\notin F}q_{ij}$. Taking $u=\ones$
gives $J_F(\alpha)\ge\sum_i\alpha_iq_i^{\rm out}(F)$, and dropping the terms
with $q_{ij}$ gives $J_F(\alpha)\le\sum_i\alpha_iq_i$.

\begin{proposition}[the restriction identity]\label{prop:restriction}
For every face $F$ and every $\alpha\in\Delta_F$, $J_F(\alpha)=J_V(\alpha)$. Hence
$J_F=J_{F'}$ on $\Delta_{F'}$ for every nonempty $F'\subseteq F$.
\end{proposition}

\begin{proof}
Fix $\alpha\in\Delta_F$ and $u\in(0,\infty)^d$. The terms of~\eqref{eq:JV} with
$i\notin F$ vanish, since $\alpha_i=0$ there. For $i\in F$ the terms with
$j\notin F$ are nonnegative and are subtracted, so
\[
 \sum_{i\in F}\alpha_i\Bigl(q_i-\sum_{j\ne i}q_{ij}\frac{u_j}{u_i}\Bigr)
 \le\sum_{i\in F}\alpha_i\Bigl(q_i-\sum_{j\in F\setminus\{i\}}q_{ij}\frac{u_j}{u_i}\Bigr)
 \le J_F(\alpha).
\]
So $J_V(\alpha)\le J_F(\alpha)$. Conversely, let $w\in(0,\infty)^F$ and let $u$
equal $w$ on $F$ and $\varepsilon$ off $F$. The terms with $j\notin F$ are then
$-q_{ij}\varepsilon/w_i$, which tend to zero as $\varepsilon\downarrow0$, and
the supremum over $w$ gives $J_V(\alpha)\ge J_F(\alpha)$. The second claim is the
first for $F$ and for $F'$.
\end{proof}

So the cost of a point of $\Delta_F$ does not depend on the face in which we
read it.

\subsection{The variational formula}\label{sec:spectral-geometry}

For a face $F$ and a field $v\in\R^F$ let
\begin{equation}\label{eq:lambdaFv}
 \lambda_F(v)=-\varrho\bigl(Q_F+\operatorname{diag}v\bigr),
\end{equation}
so that $\lambda_F(0)=\lambda_F$. If $F$ is irreducible, let $r$ and $l$ be
positive right and left Perron vectors of $Q_F+\operatorname{diag}v$ with
$l\cdot r=1$, and put $\alpha^*_F(v)=l\circ r$, so that $\alpha^*_F(0)=\alpha^*_F$.
Since $\sum_i\alpha_iv_i=v\cdot\alpha$,
\begin{equation}\label{eq:Jfield}
 J_F(\alpha)-v\cdot\alpha=\sup_{u\in(0,\infty)^F}\sum_{i\in F}\alpha_i
 \Bigl(-\frac{((Q_F+\operatorname{diag}v)u)_i}{u_i}\Bigr).
\end{equation}
So a field replaces the diagonal $q_i$ by $q_i-v_i$ and changes nothing else.

\begin{lemma}[monotonicity in the face]\label{lem:face-monotone}
If $F\subsetneq G$ are faces and $Q_G$ is irreducible, then
$\lambda_G(v_G)<\lambda_F(v_F)$ for every field $v\in\R^G$. Now let $Q$ be
irreducible. Then $\lambda_V(v)<\lambda_F(v_F)$ for every proper face $F$, and
for every fixed $v\in\R^d$ there is $\gamma(v)>0$ with
$\lambda_F(v_F)\ge\lambda_V(v)+\gamma(v)$ for all proper $F$. So the score
$\Phi_{\tau,v}(F)=|F|-1+\tau\lambda_F(v_F)$ is least at $V$ for
$\tau>(d-1)/\gamma(v)$.
\end{lemma}

The last two claims need $Q$ irreducible. For $d=2$, $q_{12}=1$ and $q_{21}=0$
we have $\lambda_V=0=\lambda_{\{2\}}$.

\begin{proof}
Choose $c$ with $B=Q_G+\operatorname{diag}v_G+cI\ge0$. Its principal submatrix on
$F$ is $Q_F+\operatorname{diag}v_F+cI$, since the diagonal of a killed generator
keeps the full exit rates. Let $\widetilde B$ be this submatrix padded with zeros
to $G\times G$. Then $0\le\widetilde B\le B$, and $\widetilde B\ne B$ because the
irreducible matrix $B$ has a positive entry in a row of $F$ and a column off $F$.
So $\varrho(\widetilde B)<\varrho(B)$~\cite[Ch.~2]{bermanplemmons1994},
and $\varrho(\widetilde B)$ is the Perron root of the submatrix on $F$, since the
padding adds only the eigenvalue $0$. This holds whether or not $F$ is
irreducible, and subtracting $c$ gives the first claim. If $Q$ is irreducible,
the first claim with $G=V$ gives the second. For fixed $v$ there are
finitely many proper faces, so $\gamma(v)$, the least of the numbers
$\lambda_F(v_F)-\lambda_V(v)$, is positive. Finally
$\Phi_{\tau,v}(F)-\Phi_{\tau,v}(V)=|F|-d+\tau(\lambda_F(v_F)-\lambda_V(v))\ge-(d-1)+\tau\gamma(v)$.
\end{proof}

\begin{proposition}[the rate on a face]\label{prop:rate}
Let $F$ be irreducible with $k=\lvert F\rvert\ge2$, and let $v\in\R^F$.
\begin{enumerate}[(a)]
\item On $\Delta_F$, $J_F(\alpha)-v\cdot\alpha\ge\lambda_F(v)$, with equality exactly at
$\alpha=\alpha^*_F(v)$. In particular $J_F\ge\lambda_F$, with equality exactly at
$\alpha^*_F=l_F\circ r_F$.
\item The rate $J_F$ is convex on $\Delta_F$ and strictly convex on $\Delta_F^\circ$. For
$\alpha\in\Delta_F^\circ$ the supremum defining $J_F(\alpha)$ is attained at a
vector $u=r_\alpha$, which is unique up to a positive factor, and $\alpha$ is
the stationary law of the generator on $F$ with off-diagonal entries
$q_{ij}r_{\alpha,j}/r_{\alpha,i}$.
\item If there is $\nu\in(0,\infty)^F$ with $\nu_iq_{ij}=\nu_jq_{ji}$ for all
$i,j\in F$, then for every $\alpha\in\Delta_F$
\[
 J_F(\alpha)=\sum_{\{i,j\}\subseteq F,\ i\ne j}\bigl(\sqrt{\alpha_iq_{ij}}-\sqrt{\alpha_jq_{ji}}\bigr)^2
 +\sum_{i\in F}\alpha_iq_i^{\rm out}(F).
\]
\item For every nonempty face $F'\subsetneq F$ and every
$\alpha\in\Delta_{F'}$, $J_F(\alpha)-v\cdot\alpha=J_{F'}(\alpha)-v\cdot\alpha\ge\lambda_{F'}(v_{F'})>\lambda_F(v)$.
\end{enumerate}
\end{proposition}

For symmetric $Q_F$, $l_F=r_F$ and $\alpha^*_F$ is the entrywise square of the
Perron vector normalized in $\ell^2$. On a face that is not reversible the
formula in (c) is only an upper bound, by its proof.

\begin{proof}
Put $B=Q_F+\operatorname{diag}v$, an irreducible Metzler matrix with the switching
graph $G_F$ and $\lambda_F(v)=-\varrho(B)$.

\emph{Setup.} With $u=e^\psi$, \eqref{eq:Jfield} reads
$J_F(\alpha)-v\cdot\alpha=\sup_\psi\Psi_\alpha(\psi)$ with
\[
\Psi_\alpha(\psi)=\sum_i\alpha_i(q_i-v_i)-\sum_{i\ne j}\alpha_iq_{ij}e^{\psi_j-\psi_i},
\]
where $i,j$ run over $F$. It is concave in $\psi$, affine in $\alpha$, and
unchanged when a constant is added to $\psi$. Let $G^\psi$ be the conservative
generator on $F$ with off-diagonal entries $q_{ij}e^{\psi_j-\psi_i}$. It has the
graph $G_F$, so it has exactly one stationary law, and it does not involve $v$.
Since $\nabla_\psi\Psi_\alpha=-\alpha^\top G^\psi$ and a differentiable concave
function attains its maximum exactly at its critical points, $\psi$ attains the
supremum if and only if $\alpha$ is the stationary law of $G^\psi$. We call this
the criterion.

\emph{A lower bound on every face.} Let $F'$ be any face, irreducible or not, and
$B'=Q_{F'}+\operatorname{diag}v_{F'}$. We show
$J_{F'}(\alpha)-v\cdot\alpha\ge-\varrho(B')$ on $\Delta_{F'}$. Let
$s>\varrho(B')$ and write $B'=B_0-cI$ with $B_0\ge0$. Then $s+c>\varrho(B_0)$, so
$u=(sI-B')^{-1}\ones=\sum_{\ell\ge0}(s+c)^{-\ell-1}B_0^\ell\ones$ converges and is
positive. Since $(sI-B')u=\ones$, we get $-(B'u)_i/u_i=-s+1/u_i>-s$ at every
state, and letting $s\downarrow\varrho(B')$ gives the claim.

\emph{The value at $\alpha^*_F(v)$.} Take $\psi=\log r$ and $R=\operatorname{diag}r$.
Then $G^\psi=R^{-1}(B-\varrho(B)I)R$, since the off-diagonal entries agree and
$Br=\varrho(B)r$. So $(l\circ r)^\top G^\psi=l^\top(B-\varrho(B)I)R=0$, and by the
criterion $\psi$ attains the supremum at $\alpha^*_F(v)$. Every ratio
$-(Br)_i/r_i$ equals $\lambda_F(v)$, so the value is $\lambda_F(v)$.

(b) Convexity holds because $J_F$ is a supremum of affine functions. Let
$\alpha\in\Delta_F^\circ$ and take $v=0$.

\emph{The supremum is attained.} Normalize $\min_i\psi_i=0$ and put
$\bar\psi=\max_i\psi_i$. Let $a$ be the least positive rate $q_{ij}$ in $F$. A
shortest path in $G_F$ from a state where $\psi=0$ to a state where
$\psi=\bar\psi$ has at most $k-1$ edges, so one of its edges $i\to j$ has
$\psi_j-\psi_i\ge\bar\psi/(k-1)$. Keeping only this edge gives
$\Psi_\alpha(\psi)\le\sum_i\alpha_iq_i-a(\min_i\alpha_i)e^{\bar\psi/(k-1)}$, which
tends to $-\infty$ as $\bar\psi\to\infty$. So the continuous function
$\Psi_\alpha$ attains its supremum on a bounded set. Write $r_\alpha=e^\psi$ for a
maximizer. The last claim of (b) is the criterion.

\emph{The maximizer is unique up to a factor.} If $\psi$ and $\psi+\delta$ both
attain the supremum, $\Psi_\alpha$ is constant on the segment between them. There
it is a constant minus a sum of the convex functions
$\alpha_iq_{ij}e^{\psi_j-\psi_i+s(\delta_j-\delta_i)}$, so each of them is affine
in $s$, and $\delta_j=\delta_i$ on every edge of $G_F$. As $G_F$ is strongly
connected, $\delta$ is constant.

\emph{Strict convexity.} Let $\alpha^0\ne\alpha^1$ lie in $\Delta_F^\circ$, put
$\alpha^s=(1-s)\alpha^0+s\alpha^1$ with $0<s<1$, and let $\psi$ attain
$J_F(\alpha^s)$. Since $\Psi_\alpha(\psi)$ is affine in $\alpha$,
$J_F(\alpha^s)=(1-s)\Psi_{\alpha^0}(\psi)+s\Psi_{\alpha^1}(\psi)\le(1-s)J_F(\alpha^0)+sJ_F(\alpha^1)$.
Equality would make $\alpha^0$ and $\alpha^1$ both stationary for $G^\psi$ by the
criterion, which is impossible.

(d) The equality is Proposition~\ref{prop:restriction}. The lower bound on $F'$
gives $J_{F'}(\alpha)-v\cdot\alpha\ge\lambda_{F'}(v_{F'})$, and
Lemma~\ref{lem:face-monotone} with $G=F$ gives $\lambda_{F'}(v_{F'})>\lambda_F(v)$.

(a) The lower bound and the value at $\alpha^*_F(v)$ are proved above. Let
$\alpha\ne\alpha^*_F(v)$. If $\alpha$ lies on the boundary of $\Delta_F$, then (d)
with $F'=\supp\alpha$ gives $J_F(\alpha)-v\cdot\alpha>\lambda_F(v)$. If
$\alpha\in\Delta_F^\circ$ and equality held, strict convexity would give a value
below $\lambda_F(v)$ at $\frac12(\alpha+\alpha^*_F(v))$.

(c) Since $q_i=\sum_{j\in F\setminus\{i\}}q_{ij}+q_i^{\rm out}(F)$, the function
inside the supremum is $\sum_i\alpha_iq_i^{\rm out}(F)$ plus, for each pair
$\{i,j\}$, the term $\alpha_iq_{ij}(1-u_j/u_i)+\alpha_jq_{ji}(1-u_i/u_j)$. By the
AM--GM inequality this term is at most
$(\sqrt{\alpha_iq_{ij}}-\sqrt{\alpha_jq_{ji}})^2$, which gives the upper bound. For
$\alpha\in\Delta_F^\circ$ the vector $u_i=\sqrt{\alpha_i/\nu_i}$ gives
$\alpha_iq_{ij}u_j/u_i=\alpha_jq_{ji}u_i/u_j$ by $\nu_iq_{ij}=\nu_jq_{ji}$, so there is
equality for every pair at once. If some $\alpha_i$ vanish, take $u_i=\varepsilon$
for these $i$ and let $\varepsilon\to0$. A pair with $\alpha_i=0<\alpha_j$ then has
the term $\alpha_jq_{ji}(1-\varepsilon/u_j)\to\alpha_jq_{ji}$, which is its bound, and
a pair with $\alpha_i=\alpha_j=0$ has the term $0$.
\end{proof}

\begin{proposition}[the variational formula]\label{prop:variational}
Let $F$ be irreducible. For every $v\in\R^F$,
\[
 \lambda_F(v)=\min_{\alpha\in\Delta_F}\bigl(J_F(\alpha)-v\cdot\alpha\bigr),
\]
and the minimum is attained only at $\alpha^*_F(v)$. Equivalently, $-\lambda_F(v)$
is the Legendre transform of $J_F$ at $v$, and $J_F$ is the Legendre transform of
$v\mapsto-\lambda_F(v)$ on $\Delta_F^\circ$. The map $v\mapsto\lambda_F(v)$ is
concave and real-analytic, and $\lambda_F(v+c\ones)=\lambda_F(v)-c$. If $Q_F$ is
reversible with respect to $\nu$, let $S_F=D_{\sqrt\nu}(-Q_F)D_{\sqrt\nu}^{-1}$, where $D_x=\operatorname{diag}x$.
Then with $\psi_i=\sqrt{\alpha_i}$,
\[
 J_F(\alpha)-v\cdot\alpha=\sum_{i\in F}(q_i-v_i)\alpha_i-\sum_{i\ne j\in F}\sqrt{q_{ij}q_{ji}}\,\psi_i\psi_j
\]
is the Rayleigh quotient of the symmetric matrix $S_F-\operatorname{diag}v$ at the
unit vector $\psi$, and $\lambda_F(v)=\lambda_{\min}(S_F-\operatorname{diag}v)$.
\end{proposition}

The formula is the finite form of the Donsker and Varadhan formula for the
principal eigenvalue~\cite{dv1975}, close to the characterization of Collatz and
Wielandt~\cite[Ch.~2]{bermanplemmons1994}.

\begin{proof}
The formula and the unique minimizer are Proposition~\ref{prop:rate}(a). As a
minimum of affine functions of $v$, $\lambda_F(v)$ is concave. It is
real-analytic because the Perron root of an irreducible Metzler matrix is a simple
eigenvalue, and the shift is clear.

\emph{Legendre duality.} The formula says
$-\lambda_F(v)=\max_\alpha(v\cdot\alpha-J_F(\alpha))$. Conversely
$J_F(\alpha)\ge v\cdot\alpha+\lambda_F(v)$ for all $\alpha$ and $v$. Let
$\alpha\in\Delta_F^\circ$, take $r_\alpha$ from Proposition~\ref{prop:rate}(b), and
put $v=q-A_Fr_\alpha/r_\alpha$, with the quotient taken entrywise. Then $Q_F+\operatorname{diag}v$ has the
off-diagonal entries of $Q_F$ and the diagonal $-A_Fr_\alpha/r_\alpha$, so
$r_\alpha>0$ is a null vector of it. Only the Perron root has a positive
eigenvector, so $\lambda_F(v)=0$. Also $J_F(\alpha)=\langle q-A_Fr_\alpha/r_\alpha,\alpha\rangle
=v\cdot\alpha$, so equality holds at this $v$.

\emph{The reversible form.} Reversibility gives
$\sqrt{\nu_i}\,q_{ij}/\sqrt{\nu_j}=\sqrt{q_{ij}q_{ji}}$, so $S_F$ is symmetric with
diagonal $q_i$ and off-diagonal entries $-\sqrt{q_{ij}q_{ji}}$. Expanding the
squares in Proposition~\ref{prop:rate}(c) gives
$J_F(\alpha)=\langle\psi,S_F\psi\rangle$, and subtracting
$v\cdot\alpha=\langle\psi,\operatorname{diag}(v)\psi\rangle$ gives the display. Every
unit vector $\psi\ge0$ arises from $\alpha=\psi\circ\psi$, and the off-diagonal
entries of $S_F-\operatorname{diag}v$ are nonpositive, so passing to the entrywise
modulus does not increase the quadratic form. Hence the minimum over $\Delta_F$ is
$\lambda_{\min}(S_F-\operatorname{diag}v)$.
\end{proof}

\subsection{The functional of the discrete chain}\label{sec:discrete-rate}

When $T$ is a fixed fraction of $N$ the counts follow the functional of the
discrete chain itself. Let $0<h<1/q_{\max}$ and $p_i=1-hq_i$, so that $I+hQ$ is
a stochastic matrix with positive diagonal. For a face $F$ and
$\alpha\in\Delta_F$ let
\[
 J^{(h)}_F(\alpha)=\frac1h\sup_{u\in(0,\infty)^F}\sum_{i\in F}\alpha_i
 \log\frac{u_i}{((I+hQ_F)u)_i},
\]
and put $J^{(0)}_F=J_F$. For $F=V$, $hJ^{(h)}_V$ is the Donsker and Varadhan rate
of the empirical measure of the chain with transition matrix
$I+hQ$~\cite{dv1975}. For $\psi\in\R^F$ put
$\theta_i(\psi)=\sum_{j\in F\setminus\{i\}}q_{ij}e^{\psi_j-\psi_i}$, and let
$L_h(y)=\int_0^1y\,ds/(1-shy)$, which is $-\log(1-hy)/h$ for $h>0$ and $y$ for
$h=0$. With $u=e^\psi$ we have $((I+hQ_F)u)_i/u_i=p_i+h\theta_i(\psi)$, so
\[
 J^{(h)}_F(\alpha)=\sup_{\psi\in\R^F}G_h(\alpha,\psi),\qquad
 G_h(\alpha,\psi)=\sum_{i\in F}\alpha_iL_h\bigl(q_i-\theta_i(\psi)\bigr),
\]
and at $h=0$ this is the supremum that defines $J_F$. Since $L_h(y)\ge y$, we
have $J^{(h)}_F\ge J_F$. We say that a function on $[0,h_*]\times U$ is
real-analytic up to $h=0$ if near each point it is the restriction of a
real-analytic function on an open set.

\begin{proposition}[the discrete functional]\label{prop:discrete-rate}
Fix $h_*$ with $h_*q_{\max}<1$.
\begin{enumerate}[(a)]
\item For every face $F$, every $0<h<1/q_{\max}$ and every $\alpha\in\Delta_F$,
$J^{(h)}_F(\alpha)=J^{(h)}_V(\alpha)$. Hence $J^{(h)}_F=J^{(h)}_{F'}$ on
$\Delta_{F'}$ for every nonempty $F'\subseteq F$.
\item Let $F$ be irreducible with $k\ge2$, and let
$(h,\alpha)\in[0,h_*]\times\Delta_F^\circ$. The supremum defining
$J^{(h)}_F(\alpha)$ is attained at exactly one $\psi=\psi^{(h)}(\alpha)$ with
$\psi_1=0$, and there the Hessian of $G_h(\alpha,\cdot)$ in
$(\psi_2,\dots,\psi_k)$ is negative definite. The maps $\psi^{(h)}(\alpha)$,
$\theta^{(h)}(\alpha)=\theta(\psi^{(h)}(\alpha))$ and $J^{(h)}_F(\alpha)$ are
real-analytic on $[0,h_*]\times\Delta_F^\circ$ up to $h=0$, and
$\theta^{(0)}(\alpha)=A_Fr_\alpha/r_\alpha$ with $r_\alpha$ from
Proposition~\ref{prop:rate}(b). The function $J^{(h)}_F$ is convex on $\Delta_F$
and strictly convex on $\Delta_F^\circ$.
\item Let $F$ be irreducible with $k\ge2$, $0\le h\le h_*$,
$\alpha\in\Delta_F^\circ$ and $\psi\in\R^F$, and put $\theta=\theta(\psi)$. Then
$\psi$ attains the supremum if and only if the vector
$(\alpha_i/(p_i+h\theta_i))_{i\in F}$ is proportional to $l\circ e^\psi$, where
$l$ is a positive left null vector of $A_F-\operatorname{diag}\theta$.
\end{enumerate}
\end{proposition}

\begin{proof}
(a) As in the proof of Proposition~\ref{prop:restriction}, the terms
$hq_{ij}u_j$ with $j\notin F$ only increase $((I+hQ)u)_i$, which gives
$J^{(h)}_V\le J^{(h)}_F$ on $\Delta_F$, and $u=\varepsilon$ off $F$ with
$\varepsilon\downarrow0$ gives the converse, since the diagonal of $I+hQ_F$ keeps
the full exit rates. The second claim is the first for $F$ and for $F'$.

(c) We have $\partial\theta_i/\partial\psi_m=q_{im}e^{\psi_m-\psi_i}$ for
$m\ne i$ and $\partial\theta_i/\partial\psi_i=-\theta_i$, and
$L_h'(y)=1/(1-hy)$. So with $w_i=\alpha_i/(p_i+h\theta_i)$ the gradient of
$G_h(\alpha,\cdot)$ is $-w^\top G^\psi$, where $G^\psi$ is the conservative
generator of the proof of Proposition~\ref{prop:rate}, with off-diagonal entries
$q_{ij}e^{\psi_j-\psi_i}$. The function $G_h(\alpha,\cdot)$ is concave by the
proof of (b) below. So $\psi$ attains the supremum if and only if $w$ is
proportional to the stationary law of $G^\psi$. With
$U=\operatorname{diag}(e^\psi)$ we have
$G^\psi=U^{-1}(A_F-\operatorname{diag}\theta)U$, whose stationary law is
proportional to $l\circ e^\psi$.

(b) \emph{Concavity.} At $h=0$ the Hessian of $-G_0$ in a direction $\delta$ is
$\sum_i\alpha_i\sum_{j\ne i}q_{ij}e^{\psi_j-\psi_i}(\delta_j-\delta_i)^2$. For
$h>0$,
\[
 -hG_h(\alpha,\psi)=\sum_{i\in F}\alpha_i
 \log\Bigl(p_i+\sum_{j\in F\setminus\{i\}}hq_{ij}e^{\psi_j-\psi_i}\Bigr).
\]
The $i$-th term is $\alpha_i$ times the logarithm of a sum of exponentials of
the affine functions $0$ and $\psi_j-\psi_i$, with the positive weights $p_i$ and
$hq_{ij}$. Its second derivative in a direction $\delta$ is $\alpha_i$ times the
variance of the values $0$ and $\delta_j-\delta_i$ under the weights
proportional to $p_i$ and $hq_{ij}e^{\psi_j-\psi_i}$. As $p_i>0$, it vanishes
only if $\alpha_i=0$ or $\delta_j=\delta_i$ on every edge $i\to j$ of $G_F$. So
for every $h\in[0,h_*]$ the function $G_h(\alpha,\cdot)$ is concave, and for
$\alpha\in\Delta_F^\circ$ its Hessian is negative definite on $\{\delta_1=0\}$,
since $G_F$ is strongly connected. This is where $h_*q_{\max}<1$ is used.

\emph{The supremum is attained.} Let $\psi_1=0$ and let $\bar\psi$ be the
largest difference of two entries of $\psi$. As in the proof of
Proposition~\ref{prop:rate}(b), some edge $i\to j$ has
$\psi_j-\psi_i\ge\bar\psi/(k-1)$, so $\theta_i(\psi)\ge ae^{\bar\psi/(k-1)}$ with
$a$ the least positive rate in $F$. Since $L_h$ is increasing and tends to
$-\infty$ at $-\infty$, the term of $i$ tends to $-\infty$ as $\bar\psi\to\infty$,
while every term is at most $\alpha_iL_h(q_i)\le q_{\max}/p_{\min}$, with
$p_{\min}=1-h_*q_{\max}$. So the supremum is attained, and by strict concavity at
exactly one $\psi$ with $\psi_1=0$.

\emph{Analyticity.} The function $L_h(y)$ is real-analytic in $(h,y)$ where
$1-hy>0$, including $h=0$ and small negative $h$, and $1-h(q_i-\theta_i)\ge p_i>0$
for $h\in[0,h_*]$. So $G$ is real-analytic near every point $(h,\alpha,\psi)$ with
$h\in[0,h_*]$. At the maximizer the gradient in $(\psi_2,\dots,\psi_k)$ vanishes
and the Hessian is invertible. The implicit function theorem gives a real-analytic
solution near $(h,\alpha)$, which is $\psi^{(h)}(\alpha)$ for $h\ge0$ by
uniqueness. Hence $\theta^{(h)}(\alpha)$ and
$J^{(h)}_F(\alpha)=G_h(\alpha,\psi^{(h)}(\alpha))$ are real-analytic up to
$h=0$. At $h=0$ the maximizer is $\log r_\alpha$, normalized.

\emph{Convexity.} As a supremum of affine functions, $J^{(h)}_F$ is convex. Let
$\alpha^0\ne\alpha^1$ lie in $\Delta_F^\circ$, let $0<s<1$, and let $\psi$ attain
$J^{(h)}_F(\alpha^s)$ at $\alpha^s=(1-s)\alpha^0+s\alpha^1$. Equality in
$J^{(h)}_F(\alpha^s)\le(1-s)J^{(h)}_F(\alpha^0)+sJ^{(h)}_F(\alpha^1)$ would make
$\psi$ attain the supremum at $\alpha^0$ and at $\alpha^1$. By (c) both would then
be proportional to $((p_i+h\theta_i)l_ie^{\psi_i})_{i\in F}$, and as probability
vectors they would be equal.
\end{proof}

For a field $v\in\R^F$, $J^{(h)}_F(\alpha)-v\cdot\alpha$ is $1/h$ times the
supremum of $\sum_i\alpha_i\log(u_i/(Mu)_i)$ with
$M=e^{h\operatorname{diag}v}(I+hQ_F)$, since $-hv_i\alpha_i=\alpha_i\log e^{-hv_i}$.
For any face $F'$ and $v\in\R^{F'}$ we put
$\lambda^{(h)}_{F'}(v)=-\frac1h\log\varrho\bigl(e^{h\operatorname{diag}v}(I+hQ_{F'})\bigr)$,
which is defined since the matrix is nonnegative with a positive diagonal.

\begin{proposition}[the discrete exit rate]\label{prop:discrete-variational}
Let $F$ be irreducible with $k\ge2$, let $0<h<1/q_{\max}$ and $v\in\R^F$, and
put $M=e^{h\operatorname{diag}v}(I+hQ_F)$. Let $r$ and $l$ be positive right and
left Perron vectors of $M$ with $l\cdot r=1$, and put
$\alpha^{*(h)}_F(v)=l\circ r$.
\begin{enumerate}[(a)]
\item On $\Delta_F$, $J^{(h)}_F(\alpha)-v\cdot\alpha\ge\lambda^{(h)}_F(v)$, with
equality only at $\alpha^{*(h)}_F(v)$, which lies in $\Delta_F^\circ$. For every
nonempty face $F'\subsetneq F$ and every $\alpha\in\Delta_{F'}$,
$J^{(h)}_F(\alpha)-v\cdot\alpha\ge\lambda^{(h)}_{F'}(v_{F'})>\lambda^{(h)}_F(v)$.
\item At $v=0$, $\varrho(M)=1-h\lambda_F$ and $\alpha^{*(h)}_F(0)=\alpha^*_F$, so
$\lambda^{(h)}_F(0)=-\frac1h\log(1-h\lambda_F)$ for every $h$.
\item Fix $h_*$ with $h_*q_{\max}<1$, and put $\lambda^{(0)}_F(v)=\lambda_F(v)$ and
$\alpha^{*(0)}_F(v)=\alpha^*_F(v)$. Then $\lambda^{(h)}_F(v)$ and
$\alpha^{*(h)}_F(v)$ are real-analytic on $[0,h_*]\times\R^F$ up to $h=0$.
\end{enumerate}
\end{proposition}

\begin{proof}
Put $B=(M-I)/h$. Its off-diagonal entries $e^{hv_i}q_{ij}$ are nonnegative and
positive exactly on the edges of $G_F$, so $B$ is an irreducible Metzler matrix.
Its entries are real-analytic in $(h,v)$ up to $h=0$, and $B=Q_F+\operatorname{diag}v$
at $h=0$. Let $\kappa=\varrho(B)$. Then $\varrho(M)=1+h\kappa$, and $M$ and $B$
have the same Perron vectors. As $\kappa$ is a simple eigenvalue of $B$, it and the
normalized Perron vectors are real-analytic in $(h,v)$ up to $h=0$, and so is
$\lambda^{(h)}_F(v)=-\kappa\int_0^1ds/(1+sh\kappa)$. This gives (c). At $v=0$ we
have $B=Q_F$ for every $h$, which gives (b).

(a) \emph{A lower bound on every face.} Let $F'\subseteq F$ be nonempty,
irreducible or not, and let $M'$ be the principal submatrix of $M$ on $F'$, which
is $e^{h\operatorname{diag}v_{F'}}(I+hQ_{F'})$ since the diagonal keeps the full
exit rates. For $\alpha\in\Delta_{F'}$, part (a) of
Proposition~\ref{prop:discrete-rate} gives
$J^{(h)}_F(\alpha)=J^{(h)}_{F'}(\alpha)$, so $J^{(h)}_F(\alpha)-v\cdot\alpha$ is
$1/h$ times the supremum of $\sum_{i\in F'}\alpha_i\log(u_i/(M'u)_i)$ over
$u\in(0,\infty)^{F'}$. Let $s>\varrho(M')$. As $M'\ge0$, the vector
$u=(sI-M')^{-1}\ones=\sum_{j\ge0}s^{-j-1}M'^j\ones$ is positive, and $M'u=su-\ones$.
Also $M'u>0$, since $M'$ has a positive diagonal. So $u_i/(M'u)_i>1/s$ for every
$i$, and letting $s\downarrow\varrho(M')$ gives
$J^{(h)}_F(\alpha)-v\cdot\alpha\ge\lambda^{(h)}_{F'}(v_{F'})$. If $F'\ne F$, then
$M'$ padded with zeros is at most $M$ and differs from it, since the irreducible
$M$ has a positive entry in a row of $F'$ and a column off $F'$. So
$\varrho(M')<\varrho(M)$~\cite[Ch.~2]{bermanplemmons1994}, which gives the second
claim. The case $F'=F$ is the lower bound.

\emph{Equality.} With $u=r$ every ratio $r_i/(Mr)_i$ is $1/\varrho(M)$. The
function $\psi\mapsto\sum_i\alpha_i(\psi_i-\log(Me^\psi)_i)$ is concave, and at
$\psi=\log r$ its derivative in $\psi_m$ is
$\alpha_m-\sum_i\alpha_iM_{im}r_m/(\varrho(M)r_i)$. For $\alpha=l\circ r$ this is
$\alpha_m-(l^\top M)_mr_m/\varrho(M)=0$. So $u=r$ attains the supremum at
$\alpha^{*(h)}_F(v)$, and equality holds there. This point lies in $\Delta_F^\circ$
since $r,l>0$.

\emph{Uniqueness.} On the boundary of $\Delta_F$ the inequality is strict by the
second claim. In $\Delta_F^\circ$, $J^{(h)}_F(\alpha)-v\cdot\alpha$ is strictly
convex by Proposition~\ref{prop:discrete-rate}(b), so it has only one minimizer.
\end{proof}
